\par\medskip\noindent\textbf{Cómo leer la salida.}
{\small\setlength{\tabcolsep}{4pt}\renewcommand{\arraystretch}{1.15}
\begin{xltabular}{\linewidth}{@{}>{\raggedright\arraybackslash}p{0.24\linewidth}>{\raggedleft\arraybackslash}p{0.15\linewidth}>{\raggedright\arraybackslash}X@{}}
\toprule Indicador & Valor & Interpretación \\ \midrule \endfirsthead
\toprule Indicador & Valor & Interpretación \\ \midrule \endhead
\bottomrule \endlastfoot
$T$ & 4615 días & Días con dato en las tres series a la vez (intersección de calendarios). \\
Matriz de correlaciones & sp-nq 0,885; sp-br 0,139; nq-br 0,107 & Correlación promedio de toda la muestra: los dos índices bursátiles se mueven casi juntos; el petróleo, casi independiente. \\
S\&P-Brent EWMA mín./máx. & $-0{,}638$ (2004) / 0,744 (2012) & La correlación EWMA ($\lambda=0{,}94$) oscila mucho: puede cambiar de signo en meses. \\
Ventana de 250 días & $-0{,}248$ a 0,544 & La ventana móvil oscila menos que la EWMA porque promedia un año completo. \\
1999--2007 / 2009--2018 & 0,006 / 0,299 & La relación acciones-petróleo apareció después de la crisis de 2008. \\
S\&P-Nasdaq EWMA & 0,487 a 0,991 & Correlación siempre alta, salvo durante la burbuja tecnológica (abril de 2000). \\
1999--2007 / 2009--2018 & 0,852 / 0,954 & También aumentó después de 2008. \\
\end{xltabular}}
